\section{Neural Network Concepts}

\begin{purposebox}
\textbf{Purpose:} Demonstrate understanding of the NN structure used for tracking control. This is \textbf{not} a deep learning NN---it is a control-theoretic neural network with specific architectural components.

\textbf{When complete:} You can describe every component of the functional block diagram and explain why the architecture is called a ``neural network.''
\end{purposebox}

%──────────────────────────────────────────────────
\subsection{Basic NN Architecture}
%──────────────────────────────────────────────────

\subsubsection*{Concept Check}

\question{Not Deep Learning}
This tracking control NN is fundamentally different from a deep learning neural network. Explain:
\begin{itemize}[nosep]
  \item In what sense is it called a ``neural network''?
  \item What plays the role of activation functions?
  \item What plays the role of weights/learning?
  \item What acts as the ``teacher'' signal (supervised learning interpretation)?
\end{itemize}

\begin{answerbox}\end{answerbox}

\question{Signum Activation Function}
Write the definition of the signum function:
\[
  \operatorname{sgn}(e) =
  \begin{cases}
    \phantom{-}?\; & \text{if } e > 0 \\
    \phantom{-}?\; & \text{if } e = 0 \\
    \phantom{-}?\; & \text{if } e < 0
  \end{cases}
\]
How does this create ``variable structure'' behavior? Compare it to smooth activation functions (sigmoid, ReLU) used in deep learning.

\begin{answerbox}\end{answerbox}

\question{Learning Rate Parameter $\beta$}
The parameter $\beta$ controls convergence speed. Explain:
\begin{itemize}[nosep]
  \item What happens if $\beta$ is too large?
  \item What happens if $\beta$ is too small?
  \item How is this analogous to learning rate in gradient descent?
\end{itemize}

\begin{answerbox}\end{answerbox}

\question{Supervised Learning Interpretation}
The reference trajectory $r(k)$ acts as the ``teacher signal.'' Explain how the tracking control loop mirrors supervised learning:
\begin{itemize}[nosep]
  \item Training data = ?
  \item Prediction = ?
  \item Error signal = ?
  \item Update rule = ?
\end{itemize}

\begin{answerbox}\end{answerbox}

%──────────────────────────────────────────────────
\subsection{Functional Block Diagram (Fig.~1)}
%──────────────────────────────────────────────────

\subsubsection*{Concept Check}

\question{Component Inventory}
The block diagram contains:
\begin{itemize}[nosep]
  \item $3n$ adders
  \item $n$ controlled switches (signum)
  \item $n$ amplifiers ($\beta$)
  \item $n$ digital differentiators ($\Delta$)
  \item $n$ blocks for $f(x)$
  \item $n \times m$ blocks for $g(x)$
  \item $n$ NN solvers (LAESS)
\end{itemize}
For a system with $n=2$ states and $m=1$ input, calculate the total number of each component.

\begin{answerbox}\end{answerbox}

\question{Key Components}
Match each component to its mathematical function:

\medskip
\begin{tabularx}{\textwidth}{l c X}
\toprule
\textbf{Component} & \textbf{Symbol} & \textbf{What it computes} \\
\midrule
Adder ($+$)          & & \\[0.8em]
Sign ($S$)           & & \\[0.8em]
Amplifier ($\beta$)  & & \\[0.8em]
Differentiator ($\Delta$) & & \\
\bottomrule
\end{tabularx}

\begin{answerbox}\end{answerbox}

\question{LAESS---Linear Algebraic Equation Solver}
The LAESS solves for $u(k)$ from the control equation. Explain:
\begin{itemize}[nosep]
  \item What equation does it solve?
  \item Why can it be implemented as a sub-network?
  \item What inputs does it need and what output does it produce?
\end{itemize}

\begin{answerbox}\end{answerbox}

\question{Signal Flow}
Trace the signal flow through the block diagram for one time step:
\begin{enumerate}[nosep]
  \item Input: current state $x(k)$ and reference $r(k)$
  \item[\textrm{2--$N$.}] \emph{Fill in each processing step.}
  \item[$N{+}1$.] Output: next state $x(k+1)$
\end{enumerate}

\begin{largeanswer}\end{largeanswer}

\question{Draw Your Own}
In the space below, sketch or describe a simplified version of the block diagram for $n=1$, $m=1$. Label every block and connection.

\begin{largeanswer}[after upper={\vspace{72pt}}]\end{largeanswer}

%──────────────────────────────────────────────────
\subsection*{Self-Assessment Checklist}
%──────────────────────────────────────────────────

\begin{itemize}[leftmargin=2em, label=$\square$]
  \item I can explain why this is called a ``neural network'' despite not being deep learning
  \item I can define the signum activation function and explain its variable structure effect
  \item I can explain the role of $\beta$ as a learning rate parameter
  \item I can list all components in the functional block diagram
  \item I can trace signal flow through the block diagram for one time step
  \item I can explain the LAESS sub-network and its purpose
  \item I can draw a simplified block diagram for a scalar system
\end{itemize}

\begin{notesbox}\end{notesbox}
